\section{Conclusion}
In this paper, we propose GNUM, a novel and general GNN-based framework with two uplift estimators for user uplift modeling. The first uplift estimators are general to different uplift scenarios, which can utilize the two groups of data together to estimate the user uplift. Specifically, when the outcome is discrete, we further propose an uplift estimator based on our designed partial labels to address the labeled data scarcity problem. Experimental results demonstrate that our proposed methods outperform state-of-the-art uplift methods under various evaluation metrics.  We also give an analysis of the relationship between the uplift and social relations. And we further demonstrate the robustness of our proposed model when labeled data is limited.  In the future, we expect to utilize more types of graphs to estimate the user uplift. 

\section{Acknowledgement}
Kun Kuang's research was supported by National Natural Science Foundation of China (62006207), Young Elite Scientists Sponsorship Program by CAST (2021QNRC001), and the Fundamental Research Funds for the Central Universities (226-2022-00142)
